\clearpage
\begin{examplebox}
\textbf{例1（已知具体双曲线：求线段比与离心率）。}

设 $H:\dfrac{x^2}{9}-\dfrac{y^2}{16}=1$，
取右焦点 $F=(5,0)$，过 $F$ 作直线平行于渐近线 $y=\frac43x$，
交右支于 $M$，交另一条渐近线于 $N$。
求 $FM/MN$ 并核对实际线段长度。

\textbf{解：}$a=3,b=4,c=5$，所以
\[
\boxed{\frac{FM}{MN}=\frac{b^2}{a^2}=\frac{16}{9}},
\qquad\boxed{e=\frac53}.
\]
用定长公式验证：
\[
\boxed{FM=\frac{8}{3}},\qquad
\boxed{MN=\frac32},\qquad
\boxed{FN=\frac{25}{6}}.
\]
因为 $8/3+3/2=25/6$，长度与 $F-M-N$ 的位置关系一致。
坐标核对：
\[
\boxed{M=\left(\frac{17}{5},-\frac{32}{15}\right)},
\qquad
\boxed{N=\left(\frac52,-\frac{10}{3}\right)}.
\]
\textbf{适用信号：}双曲线系数已知，过焦点平行渐近线，
\textbf{先直接调用定比}，不必重新联立。

\medskip
\textbf{例2（给向量比例，反求离心率与渐近线斜率）。}
已知 $\overrightarrow{FM}=3\overrightarrow{MN}$，
且构型完全符合本结论。求 $e$、$b/a$。

\textbf{解：}
\[
\lambda=3,\qquad e=\sqrt{1+3}=\boxed2,
\qquad\frac ba=\sqrt3.
\]
因此两条渐近线的斜率为 $\boxed{\pm\sqrt3}$。
若再给 $a=2$，则 $b=2\sqrt3$、$c=4$，
三段长度为 $FM=3,MN=1,FN=4$。
\end{examplebox}

\clearpage
\begin{examplebox}
\textbf{例3（题目把比例写反：先辨别分子分母）。}

同一构型，若题设为
\[
\overrightarrow{MN}=\frac9{16}\overrightarrow{FM},
\]
求双曲线离心率。

\textbf{解：}题目给的是 $MN/FM=9/16$，不是 $FM/MN$。
所以 $FM/MN=16/9$，从而
\[
\boxed{e=\sqrt{1+\frac{16}{9}}=\frac53}.
\]
不能误写成 $e=\sqrt{1+9/16}$。

\medskip
\textbf{例4（读到定长：不管 $b$ 怎样变化，$MN$ 有多长？）。}

在标准双曲线 $x^2/36-y^2/b^2=1$（$b>0$）上，
按题设作 $F,M,N$。问 $MN$ 的长度是否随 $b$ 变化？

\textbf{解：} $a=6$，故
\[
\boxed{MN=\frac a2=3}.
\]
它与 $b$ 无关。\textbf{但是} $FM=b^2/(2a)=b^2/12$
会随 $b$ 变化，焦点、直线斜率和交点坐标也会变化。

\medskip
\textbf{例5（由夹角直接求离心率）。}

设一条渐近线与 $x$ 轴的锐角为 $60^\circ$。
按题设作 $F,M,N$，求 $FM/MN$ 和 $e$。

\textbf{解：}
\[
\frac ba=\tan60^\circ=\sqrt3,
\qquad\boxed{\frac{FM}{MN}=3},
\qquad\boxed{e=\sqrt{1+3}=2}.
\]
\textbf{注意：}$60^\circ$ 是\textbf{单条渐近线与实轴}的角，
不是两条渐近线之间的某个夹角。
\end{examplebox}

\clearpage
\begin{examplebox}
\textbf{例6（迁移训练：给 $FN:MN$，不必先算 $FM$）。}

题目仍满足 C085 构型，若 $FN:MN=13:4$，求离心率。

\textbf{解：}由 $F-M-N$，$FN=FM+MN$，故
\[
\frac{FM}{MN}=\frac{FN}{MN}-1=\frac{13}{4}-1=\frac94.
\]
从而
\[
\boxed{e=\sqrt{1+\frac94}=\frac{\sqrt{13}}2}.
\]
另一条更短的路：由三段定长，
$FN/MN=(c^2/2a)/(a/2)=c^2/a^2=e^2$，
所以 $e=\sqrt{13/4}$，结果相同。

\medskip
\textbf{例7（变式：单给 $FM:FN$）。}

\begin{center}\begin{tikzpicture}[font=\small,>=stealth]
\draw[gray!45,rounded corners=3pt] (0,0) rectangle (14.2,8);
\node at (7.1,7.35) {先看整段与两部分：题目换了比值，先认分子和分母};

\draw[blue!75!black,very thick] (2,4.8)--(7,4.8);
\draw[orange!85!black,very thick] (7,4.8)--(11,4.8);
\fill (2,4.8) circle(3pt) node[above=6pt] {$F$};
\fill (7,4.8) circle(3pt) node[above=6pt] {$M$};
\fill (11,4.8) circle(3pt) node[above=6pt] {$N$};
\node[blue!75!black,below=7pt] at (4.5,4.8) {$FM:5$};
\node[orange!85!black,below=7pt] at (9,4.8) {$MN:4$};
\draw[<->] (2,3.1)--(11,3.1) node[midway,below=5pt] {$FN:9$};
\node at (7.1,1.35) {$\displaystyle\frac{FM}{MN}=\frac54,\qquad\frac{FN}{MN}=\frac94=e^2$};
\node at (7.1,.55) {比值线段图取例 7 的 $5:4:9$，不表示新的双曲线坐标。};

\end{tikzpicture}
\end{center}

图中把整段 $FN$ 看作 $9$ 份，蓝段占 $5$ 份，橙段便占 $4$ 份。故 $MN:FN=4:9$，
故 $e^2=FN/MN=9/4$，
\[
\boxed{e=\frac32},\qquad
\boxed{\lambda=\frac{FM}{MN}=\frac54}.
\]
\textbf{辨识要点：}三个常见给法可以互相换算：
\[
\frac{FM}{MN}=e^2-1,\quad
\frac{FN}{MN}=e^2,\quad
\frac{FM}{FN}=1-\frac1{e^2}.
\]
\end{examplebox}
